\section{Problem Statement}

Decentralized solutions for document verification have security problems inherent to the network on which they are deployed.
On Bitcoin, they contain privacy issues or are inefficient in cost and blockspace in solutions using protocols such as OpenTimestamps, Blockcerts, \texttt{OP\_RETURN}, among others.

Taproot applies a new, more efficient, and more private Schnorr digital signature scheme for multisignature security schemes.
It also allows committing multiple scripts into a single output using Merkleized Abstract Syntax Trees (MAST), improving costs and blockspace usage.
This combination has not been previously utilized to create a protocol for document verification.

\subsection{Research Questions}

This work aims to answer:
\begin{itemize}
    \item How can the tools applied by Taproot be leveraged to deploy a decentralized verification protocol?
    \item To what extent does using MAST improve privacy, scalability, and spatial/computational complexity?
    \item What is an appropriate methodology to add signature rotation and certificate revocation to the protocol?
\end{itemize}

\subsection{Specific Objectives}
To answer the research questions, the following objectives are set:
\begin{enumerate}[label=\alph*)]
    \item Design software that creates the Merkle structure of scripts where the integrity of the issued documents is verified.
    \item Design a software component that allows ownership verification using the digital signature on issued documents.
    \item Design a software component that allows for document signature key rotation.
    \item Design a software component that tracks revocation to verify the validity of issued documents.
    \item Deploy a Proof of Concept (PoC) on a Bitcoin test network.
    \item Measure the scalability, time/space complexity, and network costs as a function of the vBytes used in the transaction.
\end{enumerate}